% Aqui começa o capítulo de Referencial Teórico.
\chapter{Referencial Teórico}\label{chp:ReferencialTeorico}

Este capítulo apresenta os fundamentos teóricos que sustentam a abordagem de gestão de projetos adotada neste trabalho. Inicialmente, são discutidos os principais conceitos introduzidos pelo \textit{A Guide to the Project Management Body of Knowledge} em sua sétima edição, com foco em princípios e domínios de desempenho. Em seguida, são apresentados os métodos ágeis, com ênfase no framework Scrum, destacando suas características, benefícios e adequação a projetos com elevado grau de incerteza.

\section{PMBOK 7ª Edição}

A sétima edição do PMBOK representa uma mudança significativa na forma como a gestão de projetos é conceituada. Diferentemente das edições anteriores, estruturadas a partir de grupos de processos e áreas de conhecimento, o PMBOK 7 adota uma abordagem orientada por princípios e domínios de desempenho, buscando maior flexibilidade e adaptação a diferentes contextos organizacionais e tipos de projeto \cite{PMBOK7}.

Essa mudança reflete o reconhecimento de que projetos operam em ambientes cada vez mais complexos, dinâmicos e incertos, nos quais abordagens prescritivas e excessivamente padronizadas tendem a perder eficácia. Em vez de definir \textit{como} o projeto deve ser conduzido em termos de processos fixos, o PMBOK 7 estabelece \textit{diretrizes fundamentais} que orientam decisões ao longo do ciclo de vida do projeto, independentemente da abordagem adotada.

Os princípios de gerenciamento de projetos definidos pelo PMBOK 7 enfatizam aspectos como entrega de valor, engajamento efetivo dos stakeholders, adaptabilidade, liderança, colaboração, pensamento sistêmico e qualidade. Esses princípios funcionam como fundamentos normativos que devem orientar o comportamento das equipes e a tomada de decisão, permitindo que diferentes métodos, incluindo abordagens ágeis, sejam utilizados de forma coerente e alinhada aos objetivos estratégicos da organização.

Complementarmente, os domínios de desempenho descrevem áreas críticas que devem ser continuamente monitoradas e ajustadas para garantir o sucesso do projeto. Diferentemente de processos sequenciais, os domínios não representam etapas cronológicas, mas dimensões interdependentes do desempenho do projeto, como equipe, stakeholders, planejamento, entrega, incerteza e mudança. Essa visão reforça a natureza não linear dos projetos e a necessidade de gerenciamento contínuo e adaptativo.

No contexto deste trabalho, o PMBOK 7 fornece o arcabouço conceitual para avaliar a eficácia da gestão do projeto, enquanto os métodos ágeis oferecem os mecanismos práticos para operacionalizar esses princípios em um ambiente de desenvolvimento de software com requisitos emergentes.

\section{Métodos Ágeis de Gestão de Projetos}

Os métodos ágeis surgiram como resposta às limitações dos modelos tradicionais de desenvolvimento de software, especialmente em contextos caracterizados por mudanças frequentes de requisitos, alta incerteza e necessidade de entregas rápidas de valor. Fundamentados nos valores e princípios do Manifesto Ágil, esses métodos priorizam a colaboração com o cliente, a adaptação contínua e a entrega incremental de funcionalidades \cite{AgileManifesto}.

Em contraste com abordagens preditivas, que buscam definir o escopo completo do projeto antecipadamente, os métodos ágeis reconhecem que o conhecimento sobre o problema e a solução evolui ao longo do tempo. Dessa forma, o planejamento passa a ser progressivo, baseado em ciclos curtos de trabalho e revisões frequentes, permitindo incorporar feedback de maneira sistemática.

Diversos frameworks e práticas ágeis foram desenvolvidos ao longo do tempo, entre eles Scrum, Kanban e abordagens híbridas. Cada um desses métodos apresenta características específicas, mas compartilham princípios comuns, como foco em valor, transparência, inspeção e adaptação.

\section{Scrum como Framework Ágil}

Scrum é um framework ágil amplamente utilizado para o gerenciamento e desenvolvimento de produtos complexos. Sua estrutura baseia-se em ciclos iterativos e incrementais chamados de sprints, nos quais uma equipe multifuncional trabalha para entregar incrementos potencialmente utilizáveis do produto \cite{ScrumGuide}.

O framework define papéis, eventos e artefatos com o objetivo de promover transparência, inspeção e adaptação contínuas. O Product Backlog representa a lista priorizada de necessidades do produto, enquanto o Sprint Backlog detalha o trabalho selecionado para o ciclo corrente. Eventos como a Sprint Review e a Sprint Retrospective permitem avaliar os resultados obtidos e ajustar o processo de trabalho conforme necessário.

Uma característica central do Scrum é a ênfase na auto-organização da equipe e na colaboração constante com os stakeholders. Em vez de um controle centralizado e prescritivo, o framework estimula a responsabilidade compartilhada e a tomada de decisão baseada em evidências empíricas obtidas ao longo da execução do projeto.

No contexto do sistema de gestão de casos sociais proposto neste trabalho, o Scrum mostra-se adequado por permitir a incorporação gradual de funcionalidades, a validação contínua com usuários e gestores, e a adaptação do backlog frente a mudanças regulatórias, operacionais ou de contexto social. Além disso, o uso de métricas ágeis contribui para o aumento da previsibilidade e do controle do projeto, em alinhamento com os domínios de desempenho do PMBOK 7.

\section{Integração entre PMBOK 7 e Métodos Ágeis}

A integração entre o PMBOK 7ª edição e métodos ágeis não deve ser compreendida como uma sobreposição de modelos, mas como uma relação complementar. Enquanto o PMBOK 7 estabelece os princípios orientadores e os resultados esperados em termos de desempenho, os métodos ágeis, como o Scrum, oferecem práticas concretas para alcançar esses resultados em ambientes dinâmicos.

Essa integração permite realizar o \textit{tailoring} da abordagem de gestão do projeto, selecionando e adaptando práticas de acordo com o contexto, os riscos e as características da organização. Assim, torna-se possível combinar governança, alinhamento estratégico e controle com flexibilidade, aprendizado contínuo e foco em valor.

Dessa forma, o referencial teórico adotado neste trabalho sustenta a escolha de uma abordagem ágil orientada por princípios, adequada à complexidade do projeto de gestão de casos sociais e alinhada às diretrizes contemporâneas da gestão de projetos.

